We built a guarded, reproducible four-city benchmark of next-day reported
311 demand and coupled it to a controlled, explicitly simulated
capacity-allocation evaluation with frozen pre-test environments and
validation-only selection. The corrected evidence supports a precise set
of claims: public calendar signals improve next-day forecasts in every
system and weather helps heterogeneously; temporally valid pooling and
transfer mostly hurt; within the simulation, \emph{how} a forecast
enters the allocation rule can matter as much as which forecaster
produced it --- a degenerate point forecast leaves the expected-value
allocation objective non-identified, so its realized loss is determined
by the secondary tie-break rule, and either a proportional secondary rule
or the quantile-interpolated predictive distribution of a single quantile
model recovers a simple proportional policy's performance under the
tested regimes and neutral objective; and decision-aware validation
selection mostly coincides with accuracy-based selection and did not
increase test loss in the pre-specified grid.
The broader lesson for forecast-driven public operations research is
methodological: evaluation environments must be frozen before selection,
selection must never touch test outcomes, and the interface between a
forecast and a policy deserves the same scrutiny as the forecast itself.
Natural continuations include richer service-time modeling where
case-closure data permit, spatial disaggregation where geographies allow
honest comparison, and extension beyond U.S. systems as comparable public
data emerge.
